%% STORM-PhysNet -- IEEE CONFERENCE (target ~6 pages)
\documentclass[conference]{IEEEtran}
\usepackage{amsmath,amssymb,amsfonts}
\usepackage{graphicx}
\usepackage{booktabs}
\usepackage{multirow}
\usepackage{cite}
\usepackage{url}
\usepackage{xcolor}
%\usepackage{balance}

\title{STORM-PhysNet: A Multi-Horizon Transformer for Geostationary Relativistic Electron Flux Forecasting with Physics-Inspired Components and Cross-Satellite Transfer}

\author{
\IEEEauthorblockN{Samarth BN}
\IEEEauthorblockA{Department of Electronics and Communication Engineering\\
RV College of Engineering, Bengaluru, India\\
Email: samarthbn.ec25@rvce.edu.in}
}

\begin{document}
\maketitle

%% --------------------------------------------------------------------
\begin{abstract}
Energetic electrons above 2\,MeV at geostationary Earth orbit (GEO)
are a persistent internal-charging hazard for satellites. We present
STORM-PhysNet, a multi-horizon forecast model that couples a standard
Transformer encoder over a 72-hour window with a learnable L1--Earth
propagation delay, a \(B_z\)-conditioned physics gate, and residual
heads for simultaneous 1-hour, 6-hour, and 12-hour log-flux
forecasts.

Relative to a width/depth/head-matched Transformer
(\(d_{\mathrm{model}}=128\), two layers, four heads; not FFN-width
matched and without STORM's persistence residual), hereafter the
matched Transformer, STORM-Bz retains a small
1\,h PE$_{\mathrm{clim}}$ edge (0.986 vs 0.980) and a large advantage versus
persistence (PE$_{\mathrm{pers}} \approx 0.31$ vs $\approx -0.01$). At 6\,h and 12\,h, single-model
means are close (STORM-Bz 0.900 / 0.854 vs matched 0.895 / 0.845).
Seed-bagged STORM-Bz reaches PE$_{6\mathrm{h}} = 0.910$ and PE$_{12\mathrm{h}} = 0.870$;
bagged matched Transformer is competitive (0.908 / 0.861).
A default Transformer baseline ($d_{\mathrm{model}}=64$, three layers; not
capacity-matched) scores PE$_{1\mathrm{h}} = 0.978$ and is reported only as a
secondary reference. Controlled ablations (No-Delay, No-Gate,
No-Physics) show that the short-horizon gain is a property of the
trained package as a whole rather than of any individual physics
module at inference.

Under synthetic solar-wind input noise and targeted channel dropout,
STORM degrades more slowly than the Transformer, especially at
1-hour. On GSAT-19 GRASP, fine-tuning raises 6\,h PE from
0.740 to 0.841 (paired \(p = 3.55\times10^{-12}\)) and 12\,h PE from
0.567 to 0.762 (\(p = 1.66\times10^{-13}\)) on a $\sim$14--15 month GRASP window (fifteen seeds; paired tests on seed-level PE).
Bootstrap confidence intervals, paired tests, and evaluation scripts
are released so that later work can reuse the same splits and PE
definitions.
\end{abstract}

\begin{IEEEkeywords}
Geostationary orbit, relativistic electrons, space weather, Transformer,
physics-inspired, transfer learning, prediction efficiency, GRASP
\end{IEEEkeywords}

%% ====================================================================
\section{Introduction}
%% ====================================================================
Geostationary satellites sit inside the outer radiation belt, where
relativistic electrons (\(E>2\,\mathrm{MeV}\)) can jump by several orders of
magnitude in a few hours when a geomagnetic storm develops. Those bursts drive
internal charging and can upset or damage onboard electronics~\cite{b1,b2}.
For operators of GEO assets at Indian longitudes---ISRO broadcasting,
navigation, and remote-sensing satellites among them---a forecast that is
usable at the local meridian, not only at GOES longitudes, is increasingly
useful.

Most published skill scores still come from GOES-centered data sets. When a
new instrument such as GRASP on GSAT-19 comes online, there is often only a
short calibrated record, so operators need a practical path from a data-rich
satellite to a data-scarce one without retraining everything from scratch.

Forecast models have progressed from linear filters~\cite{b3,b4}
through LSTM-based deep-learning approaches~\cite{b5,b6,b7}
to recent Transformer-based architectures reporting daily-fluence PE of
0.85--0.92~\cite{b8,b9,b10}. We do not include REFM as a numeric baseline because it targets daily integral
fluence rather than hourly log-flux; a direct PE comparison would be
misleading~\cite{b4}. Despite these advances,
two gaps remain: (i) the short-horizon (\(\sim\)1\,h) behavior of
Transformer architectures relative to physics-inspired variants is poorly
characterized across seeds, and (ii) no systematic study has assessed
transfer strategies for Indian-longitude GEO assets such as GSAT-19.

The design augments a standard Transformer with constrained delay and
gating modules, trained end-to-end under matched chronological splits
and fifteen seeds. Primary comparisons use an architecture-matched
Transformer; a default-width Transformer and LSTM are reported as
references. Ablations in Section~\ref{sec:results} show that no single
module fully accounts for the short-horizon PE$_{\mathrm{clim}}$ edge.

This manuscript focuses on the multi-seed GOES benchmark, the
short-horizon reliability gap, the main ablations, and the GRASP transfer
recipe that is most relevant to Indian-longitude operators.

%% ====================================================================
\section{Related Work}
%% ====================================================================
Camporeale~\cite{b11} summarizes common failure modes when ML is
applied to space weather (leaky splits, mismatched metrics, overstated skill).
Prediction efficiency (PE) remains the standard skill score~\cite{b12}, but care must be taken to distinguish between climatological and persistence baselines.

On the modeling side, work moved from linear filters and empirical solar-wind
drivers~\cite{b3,b4,b13} to LSTM and hybrid deep
models~\cite{b5,b6,b7}, and more recently to
Transformer-style networks for GEO flux and fluence~\cite{b8,b9,b10}.
Among recent Transformer studies, Sun et al.~\cite{b10} provide the closest multi-horizon GEO baseline; our evaluation differs by using a matched fifteen-seed protocol, explicit delay/gate ablations, and a documented GOES$\rightarrow$GRASP transfer path.
Those papers mainly report daily or multi-hour skill on GOES-like targets.
Short-lead (about 1\,h) skill is less often reported with multi-seed error bars,
and cross-satellite transfer to a short Indian-longitude record is rarely
benchmarked under the same horizons and persistence definition. Physics-inspired
constraints, when used, are often loss penalties rather than an explicit
delay-plus-gate structure trained end-to-end with multi-horizon residual heads.
We keep a standard temporal Transformer backbone so that comparisons
isolate the effect of adding constrained delay and gating under a shared
encoder family. Ablations in
Section~\ref{sec:results} show that no single module fully accounts for the short-horizon
PE\(_{\mathrm{clim}}\) edge.

Solar-wind timing is another recurring issue. OMNI applies a
phase-front correction from L1, yet residual window-to-window lags remain
a known uncertainty in solar-wind--magnetosphere coupling studies. We treat that lag as a dynamically
predicted parameter rather than a fixed offset. Meanwhile, cross-platform radiation-belt products
exist for other fleets~\cite{b15}; the specific GOES\(\rightarrow\)GRASP
path at \(\approx 82^\circ\)E is the transfer setting studied here.

\begin{figure*}[!t]
\centering
\includegraphics[width=0.95\textwidth]{figures/fig_system_architecture.png}
\caption{STORM-PhysNet. OMNI solar wind is delay-aligned ($\tau\in[0.5,1.5]$~h) and
concatenated with GOES log-flux ($T=72$~h), projected to $d_{\mathrm{model}}=128$,
and encoded by a 2-layer, 4-head Transformer. A $B_z$-conditioned gate modulates
the last hidden state; three residual heads add a persistence skip at 1~h, 6~h,
and 12~h.}
\label{fig:arch}
\end{figure*}

%% ====================================================================
\section{Data and Metrics}
%% ====================================================================
\textbf{GOES-15.} We use \(E>2\,\mathrm{MeV}\) integral electron flux from a
chronological working sample spanning approximately March 2010 to February 2020
within the GOES-15 archive. One-minute samples are aggregated to hourly means;
the within-hour standard deviation is kept as an input feature. Flux is
\(\log_{10}\)-transformed. The dataloader uses sixteen solar-wind /
geomagnetic channels
(\(V_{sw}\), \(B_z\), \(B_t\), density, \(P_{\mathrm{dyn}}\),
\(B_z\)-south duration, storm-onset hours, rolling \(V_{sw}\) and \(B_z\)
moments, \(B_z\times V_{sw}\), Kp, Dst, and within-hour log-flux std)
plus log-flux history (17 inputs total after concatenation). Forecast targets use integer hourly leads of 1\,h, 6\,h, and 12\,h on the hourly GOES series (config/dataloader: $[1,6,12]$). A 72\,h lookback is used at
each forecast origin. All derived channels and indices in the 72\,h window are values
already available at the forecast origin; no future $K_p$ or
$D_{\mathrm{st}}$ is used as an input. The released tables use a fixed chronological cut of
\(63\,394\) / \(7\,924\) / \(15\,850\) hours (no shuffling).
The preprocessor's generic 70/15/15 helper is not that cut.
Test PE is computed once and never used for selection.

\textbf{GRASP} (GSAT-19, Indian longitude \(\approx82^\circ\)E) provides
$\approx$14--15 months of GEO radiation measurements from July 2017 to September 2018 in
the declining phase of Solar Cycle 24. After alignment, GRASP supplies fewer
co-located input channels than the GOES pipeline; this mismatch is a main
reason zero-shot transfer fails and motivates input-projection adaptation.
Storm windows on GRASP are defined by Dst \(\le -50\,\mathrm{nT}\)
(118/1781 test hours, 6.6\%).

\textbf{Skill score.} We report two PE definitions so that results remain comparable both to climatological skill and to persistence-based operational baselines~\cite{b12}:
\begin{align}
\mathrm{PE}_{\mathrm{clim}} &= 1 - \frac{\mathrm{MSE}(\hat y,y)}{\mathrm{Var}(y)}, \\
\mathrm{PE}_{\mathrm{pers}} &= 1 - \frac{\mathrm{MSE}(\hat y,y)}{\mathrm{MSE}(y^{\mathrm{pers}},y)}.
\end{align}
Unless stated otherwise, tables use \(\mathrm{PE}_{\mathrm{clim}}\) with fifteen seeds and bootstrap 95\% confidence intervals. The fifteen seeds measure sensitivity to random initialization on a fixed chronological split; they do not constitute independent trials across different solar-cycle phases or data periods. \(\mathrm{PE}_{\mathrm{pers}}\) is emphasized at the 1-hour horizon. Storm windows use Dst \(\le -50\,\mathrm{nT}\).

%% ====================================================================
\section{STORM-PhysNet Architecture}
%% ====================================================================

\subsection{Adaptive Propagation Delay}
Solar-wind transit from L1 to the magnetopause is typically on the
order of \(0.5\)--\(1.5\,\mathrm{h}\) in OMNI phase-front-corrected
data~\cite{b14}. Rather than a fixed offset, STORM-PhysNet predicts a per-window delay \(\tau = f_\theta(X_{\mathrm{sw}})\) from a short conditioning network on the upstream solar-wind channels. The delay module is constrained to \([0.5, 1.5]\,\mathrm{h}\), a band chosen to cover the transit range while still allowing limited storm-time deviation without breaking causality. Differentiable linear interpolation then time-shifts the solar-wind features. We do not claim a unique recovered constant of 1.09\,h from every window.

\subsection{\(B_z\)-Conditioned Physics Gate}
Sustained southward IMF \(B_z\) triggers dayside reconnection and ring-current
injection~\cite{b2}. A lightweight gate maps a short vector of
southward-\(B_z\) statistics through an MLP and a clamped sigmoid to a
per-channel modulation vector \(\mathbf{g}\in[g_{\min},1]^{d_{\mathrm{model}}}\), which is multiplied element-wise with the encoder representation
before the forecast heads. Gate activations differ between storm and quiet windows
(Mann--Whitney \(p = 1.2\times10^{-14}\), see Section~\ref{sec:results}); the absolute range
is narrow.

\subsection{Transformer Backbone and Multi-Horizon Heads}
We use a standard Transformer encoder~\cite{b16}: each of the \(T=72\)
time steps is a token; self-attention runs across time; the final hidden state
feeds the gate and three residual heads for the 1\,h, 6\,h, and 12\,h
forecasts. Measured parameter counts are approximately $3.43\times10^{5}$ for
STORM-Bz and $8.45\times10^{5}$ for the default Transformer baseline
($d_{\mathrm{model}}=64$, 3 layers, 4 heads). An architecture-matched Transformer
control ($d_{\mathrm{model}}=128$, 2 layers, 4 heads; $\approx1.19\times10^{6}$
parameters) reaches mean PE$_{1\mathrm{h}}=0.980$ and PE$_{6\mathrm{h}}=0.895$
(fifteen seeds); bagging yields $0.984$ / $0.908$ / $0.861$. Matching width/depth
shrinks the short-horizon gap versus STORM-Bz and improves multi-horizon PE under
bagging. The default baseline is not architecture-matched; single-model PE differences
should be interpreted with this disparity in mind.

\subsection{Physics Regularization and Training}
Training minimises a composite objective: a primary multi-horizon residual loss (beta-NLL with a \(2.5\times\) under-prediction penalty during storm periods), plus auxiliary Dst/Kp and storm-onset terms and small regularisers (monotonicity, temporal smoothness, delay, variance calibration, and an Akasofu-\(\varepsilon\) coupling consistency term). Because $B_y$ is absent, the Akasofu-$\varepsilon$ term approximates
the clock angle from $B_z$ and $B_t$ and is a relative regularizer,
not a calibrated coupling estimate. Horizon weights emphasise the short horizon. Coefficient values follow the released configuration. We did not ablate every auxiliary term individually; reported PE reflects the full composite objective used in training. All primary models use AdamW, cosine decay, early stopping on the 12\,h validation MSE, and seeds \(\{42,\ldots,56\}\).

\subsection{Baselines}
The primary control is an architecture-matched Transformer
($d_{\mathrm{model}}=128$, two layers, four heads;
$\approx1.19\times10^{6}$ parameters). A default-width Vanilla Transformer
($d_{\mathrm{model}}=64$, three layers) and an LSTM are reported as
secondary references.




%% ====================================================================
\section{Results}
\label{sec:results}
%% ====================================================================

\subsection{GOES Multi-Seed Performance}
Table~\ref{tab:main} summarizes fifteen-seed test results.
The primary comparison is the architecture-matched Transformer
(PE$_{1\mathrm{h}}=0.980$, PE$_{6\mathrm{h}}=0.895$).
STORM-Bz retains a small 1\,h PE$_{\mathrm{clim}}$ edge ($0.986$)
and a large PE$_{\mathrm{pers}}$ advantage ($+0.31$ vs $-0.01$).
A default-width Transformer scores $0.978$ at 1\,h and is secondary.
Seed-bagged STORM reaches PE$_{6\mathrm{h}}=0.910$ and PE$_{12\mathrm{h}}=0.870$.
Validation selects a mixing weight \(\alpha^*\) per seed (typically 0.4--0.6). We also report a fixed \(\alpha=0.3\) diagnostic. Under selected \(\alpha^*\), linear-ensemble test PE\(_{6\mathrm{h}} \approx 0.906\). We do not claim online \(\alpha\) adaptation.
Bagging fifteen STORM seeds (\(0.910\) at 6\,h) is stronger than the hybrid ensemble and provides a similar multi-horizon profile without requiring the Transformer at inference. We recommend seed-bagging for deployment because it needs only one
architecture at inference and matches or exceeds the hybrid
$\alpha^\ast$ ensemble.

\begin{table}[!htbp]
\centering
\caption{GOES test PE (fifteen seeds). Bold: highest value among the rows
shown in that column. PE is $\mathrm{PE}_{\mathrm{clim}}$ except
the last column.}
\label{tab:main}
\setlength{\tabcolsep}{3.2pt}
\begin{tabular}{lccccc}
\toprule
System & PE\(_{1\mathrm{h}}\) & PE\(_{6\mathrm{h}}\) & PE\(_{12\mathrm{h}}\) & PE\(_{\mathrm{st},6\mathrm{h}}\) & PE\(_{\mathrm{pers},1\mathrm{h}}\) \\
\midrule
LSTM                         & \(0.955\) & \(0.881\) & \(0.840\) & \(0.788\) & \(-1.19\) \\
Transformer (default)        & \(0.978\) & \(0.895\) & \(0.845\) & \(0.797\) & \(-0.08\) \\
STORM-Bz                     & \(0.986\) & \(0.900\) & \(0.854\) & \(0.812\) & \(+0.31\) \\
STORM (No-Delay)             & \(\mathbf{0.987}\) & \(0.902\) & \(0.855\) & \(0.809\) & \(\mathbf{+0.36}\) \\
STORM (No-Physics)           & \(0.986\) & \(0.899\) & \(0.850\) & \(0.807\) & \(+0.31\) \\
STORM (No-Gate)              & \(0.986\) & \(0.900\) & \(0.856\) & \(0.819\) & \(+0.32\) \\
Transformer matched          & \(0.980\) & \(0.895\) & \(0.845\) & \(0.809\) & \(-0.01\) \\
\midrule
STORM-Bz bagged              & \(\mathbf{0.987}\) & \(\mathbf{0.910}\) & \(\mathbf{0.870}\) & \(\mathbf{0.836}\) & --- \\
TF matched bagged            & \(0.984\) & \(0.908\) & \(0.861\) & \(0.831\) & --- \\
TF default bagged           & \(0.984\) & \(0.907\) & \(0.861\) & \(0.816\) & --- \\
\bottomrule
\end{tabular}

\vspace{1ex}
{\footnotesize \raggedright Means over seeds 42--56.
Bagged rows average the fifteen predictions before scoring. PE is \(\mathrm{PE}_{\mathrm{clim}}\).\par}
\end{table}

\begin{figure}[t]
\centering
\includegraphics[width=\columnwidth]{figures/fig_horizon_pe.png}
\caption{Per-horizon PE against climatology (mean\(\pm\)std, fifteen seeds).
STORM-Bz leads at 1\,h; at 6\,h and 12\,h the single models are close,
while the ensemble improves multi-horizon PE (see Table~\ref{tab:main}).}
\label{fig:horizon}
\end{figure}

\subsection{Ablation Analysis}
In all STORM variants, physics-loss terms are restricted to the
short-horizon head by default, as early testing showed that enforcing
these constraints on the 6\,h and 12\,h heads did not improve long-horizon PE.

A No-Gate run (identity modulation) retains 1\,h PE near \(0.986\), so the short-horizon PE\(_{\mathrm{clim}}\) advantage over
the Transformer is not removed by turning the gate off alone.

On storm windows, single-model PE\(_{\mathrm{st},6\mathrm{h}}\) values
are close: STORM-Bz \(0.812\), LSTM \(0.788\), Transformer \(0.797\),
and No-Gate at \(0.819\). We do not claim any one physics
module is required for storm-time skill. The clearest gains on storm
subsets remain post-hoc aggregation: the ensemble under selected \(\alpha^*\)
reaches \(0.822\) and bagged STORM reaches \(0.836\).

A single-seed diagnostic stratified by Dst intensity (full working
sample) shows PE\(_{6\mathrm{h}}\) remaining above 0.70 even in the
strong-storm bin (\(\mathrm{Dst}\le-100\,\mathrm{nT}\)). Skill is not
confined to quiet, persistence-dominated intervals.

The ablations together indicate that the short-horizon gain is produced by the training protocol itself---storm-biased sampling, the storm-weighted loss term, and residual heads with a skip from the last observed flux---rather than by any single runtime module. The delay and gate act as structured inductive biases during training; once the model is trained, removing either module at inference does not erase the 1\,h PE\(_{\mathrm{clim}}\) edge over the Transformer. We therefore present STORM-PhysNet primarily as a carefully engineered training and multi-horizon recipe that includes interpretable physics-inspired components, rather than as a demonstration that the delay or gate modules are individually necessary for the observed gains. The FFN-matched control in Limitations shows that adding a persistence
residual, not the delay or gate, is what closes the 1\,h gap to STORM.

\subsection{Input Noise and Missing Channels}
\begin{figure*}[t]
\centering
\includegraphics[width=7.16in]{figures/noise_robustness.png}
\caption{Prediction efficiency (1\,h and 6\,h) under synthetic Gaussian noise (\(\sigma\)) injected into standardized solar-wind inputs.}
\label{fig:noise}
\end{figure*}
Under synthetic input noise (Fig.~\ref{fig:noise}), the architecture-matched Transformer's 6\,h PE drops to \(0.825\pm 0.010\) at \(\sigma=1.0\), while STORM maintains \(0.848\pm 0.008\). These noise-probe values are not a re-statement of Table~\ref{tab:main}; \(\sigma=0\) under the probe protocol is not the official seed-mean PE.

\subsection{Interpretability}
Predicted delays are constrained to \([0.5, 1.5]\,\mathrm{h}\). Across the full test set the delay distribution shows a non-trivial fraction of windows concentrating near the upper bound (\(1.5\,\mathrm{h}\)), with the remainder spread across the interior of the constraint band. This saturation at the upper bound indicates the box constraint is active; we therefore treat the histogram as a constrained aggregate rather than evidence of an unconstrained recovered transit time. Gate activations differ between storm and quiet windows (Mann--Whitney \(p = 1.2\times10^{-14}\)); the absolute range remains narrow, so the shift is statistically clear but modest in magnitude. Fig.~\ref{fig:case_enhancement} shows an illustrative enhancement event in which the model tracks the multi-day flux build-up and recovery more closely than persistence.
Permutation importance ranks \(B_z\) first among inputs, consistent with the
gate design. Residual histograms on the test set widen under storm conditions for both
models; STORM reduces the short-horizon bias relative to persistence but
does not eliminate storm-time residual variance.

\begin{figure*}[t]
\centering
\includegraphics[width=7.16in]{figures/fig_case_study_timeseries.png}
\caption{Illustrative enhancement event case study. The STORM-Bz model captures the multi-day flux build-up and recovery better than the persistence baseline.}
\label{fig:case_enhancement}
\end{figure*}

\subsection{GRASP Transfer}
On GSAT-19 GRASP (fifteen seeds), zero-shot GOES models already achieve
short-horizon PE near \(0.96\), but 6\,h and 12\,h PE fall to
\(0.740\) and \(0.567\). Fine-tuning on GRASP (freezing the delay module, \(B_z\) gate, and Transformer encoder
and updating only the forecast heads) raises 6\,h PE from
\(0.740 \pm 0.020\) to \(0.841 \pm 0.017\)
(paired \textit{t}-test across fifteen seeds,
\(p = 3.55\times10^{-12}\), Cohen's \(d = 5.60\))
and 12\,h PE from \(0.567 \pm 0.033\) to \(0.762 \pm 0.020\)
(\(p = 1.66\times10^{-13}\), \(d = 7.00\)). These tests measure seed-to-seed consistency on one GRASP interval,
not independent solar-cycle trials; all fifteen seeds improve after
fine-tuning. Absolute GRASP PE remains below
the GOES test scores, as expected under sensor and longitude shift; the
practical point is the large mid- and long-horizon recovery after
adaptation rather than a claim of parity with GOES. A similar adaptation profile is observed for the architecture-matched Transformer baseline, which improves its 6\,h PE from \(0.797 \pm 0.014\) (zero-shot) to \(0.878 \pm 0.006\) after fine-tuning.

\begin{table}[t]
\centering
\caption{GRASP (GSAT-19) transfer results (fifteen seeds). Bold: best per horizon.}
\label{tab:grasp}
\begin{tabular}{lccc}
\toprule
Horizon & Zero-shot & Fine-tuned & Gain \\
\midrule
1 h  & 0.964 $\pm$ 0.001 & 0.968 $\pm$ 0.001 & +0.005 \\
6 h  & 0.740 $\pm$ 0.020 & \textbf{0.841 $\pm$ 0.017} & \textbf{+0.102} \\
12 h & 0.567 $\pm$ 0.033 & \textbf{0.762 $\pm$ 0.020} & \textbf{+0.195} \\
\bottomrule
\end{tabular}

\vspace{1ex}
{\footnotesize \raggedright Values are means $\pm$ standard deviation over seeds 42--56. Paired tests across fifteen seeds are reported in the text.\par}
\end{table}

\begin{figure}[t]
\centering
\includegraphics[width=\columnwidth]{figures/fig_grasp_domain_gap.png}
\caption{GRASP domain-gap summary (fifteen seeds): zero-shot GOES models versus
fine-tuning on GSAT-19. Short-horizon PE transfers more readily; 6\,h and 12\,h
PE improve substantially after adaptation.}
\label{fig:grasp}
\end{figure}

%% ====================================================================
%% ====================================================================
\section{Limitations and Conclusion}

\textbf{Resolution.} All reported leads are integer hours on hourly data; sub-hour (e.g.\ 45\,min) supervision would require higher-cadence flux labels.
\textbf{Data scope.} Main GOES tables use fifteen seeds with bootstrap intervals; secondary diagnostics (individual case windows, residual histograms) are illustrative. GRASP covers a limited declining-phase window; performance during solar maximum is unknown.

\textbf{Comparability.} PE values are not interchangeable with daily-fluence scores reported elsewhere without matching cadence, energy channel, and baseline~\cite{b11,b12}. The default Transformer baseline is not architecture-matched to STORM; a matched control ($d_{\mathrm{model}}=128$, two layers) shows that encoder width/depth and bagging account for much of the multi-horizon PE movement. The architecture-matched Transformer copies STORM's $d_{\mathrm{model}}$,
layer count, and head count. It does not copy the feed-forward width
($d_{\mathrm{ff}}=256$ in STORM versus PyTorch's default
$\mathrm{dim\_feedforward}=2048$). Neither baseline uses STORM's
residual-over-persistence heads. A follow-on FFN-matched control ($d_{\mathrm{ff}}=256$) remains below
STORM at 1\,h (PE$_{1\mathrm{h}}=0.976$ versus $0.986$); adding a
persistence residual closes that gap. Details are in the released
evaluation CSVs.

\textbf{Ensemble.} Ensemble skill uses a validation-selected \(\alpha^*\); the test-set \(\alpha\) curve and \(\alpha=0.3\) are diagnostic only. Online \(\alpha\) adaptation is not claimed.

\textbf{Statistical inference.} The reported p-values are uncorrected for multiple comparisons across horizons and ablations; treat individual significance claims as descriptive rather than confirmatory. All architecture comparisons are internal (STORM versus our own Transformer and LSTM implementations on the same splits). We do not claim to surpass published PE numbers from prior work, which use different cadences, energy channels, and evaluation periods.

Relative to the architecture-matched Transformer, the clearest single-model gain is short-horizon
reliability (PE$_{1\mathrm{h}}$ and PE$_{\mathrm{pers}}$).
Among the bagged rows in Table~\ref{tab:main}, STORM-Bz reaches
$\mathrm{PE}_{6\mathrm{h}}=0.910$ and $\mathrm{PE}_{12\mathrm{h}}=0.870$.
No-Delay and No-Gate bagged scores in the released CSV differ by
less than $0.002$; we do not treat STORM-Bz as uniquely best under
bagging. Seed-bagging is still the practical choice: one architecture
at inference and no validation-time $\alpha^\ast$.
Robustness checks and GRASP fine-tuning support operational motivation beyond a pure leaderboard PE number; for instance, the architecture-matched baseline remains near persistence at 1\,h (PE$_{\mathrm{pers}}\approx-0.01$), whereas STORM-Bz is clearly positive (PE$_{\mathrm{pers}}\approx+0.31$).



\section*{Data Availability}
The code, configuration files, trained checkpoints, evaluation scripts,
result CSVs, paper figures, and the chronological split protocol are
available at
\url{https://github.com/bnsama29-cloud/STORM-PhysNet-L1}. Note that the GRASP fine-tuning run was executed externally, and its frozen outputs are provided in the repository.
The GOES-15 EPEAD dataset is publicly available from NOAA NCEI
(\url{https://www.ngdc.noaa.gov/stp/satellite/goes/}).
OMNI is available via NASA OMNIWeb
(\url{https://omniweb.gsfc.nasa.gov/}).
GSAT-19 GRASP is accessible via ISSDC
(\url{https://www.issdc.gov.in/}). 

\section*{Acknowledgments}
The authors thank the GOES, OMNI, and GRASP data providers.

\begin{thebibliography}{99}
\bibitem{b1}
D.~N. Baker et al., ``Space weather effects in the Earth's radiation belts,''
\emph{Space Sci.\ Rev.}, vol.~214, p.~17, 2018.
\bibitem{b2}
R.~B. Horne et al., ``Wave acceleration of electrons in the Van Allen
radiation belts,'' \emph{Nature}, vol.~437, pp.~227--230, 2005.
\bibitem{b3}
D.~N. Baker et al., ``Linear prediction filter analysis of relativistic
electron properties at 6.6 \(R_E\),'' \emph{J. Geophys. Res.}, vol.~95,
pp.~15133--15140, 1990.
\bibitem{b4}
K.~A. Perry et al., ``Comparing geosynchronous relativistic electron
prediction models,'' \emph{Space Weather}, vol.~8, S12003, 2010.
\bibitem{b5}
L.~Wei et al., ``Quantitative prediction of high-energy electron integral flux
at GEO based on deep learning,'' \emph{Space Weather}, vol.~16,
pp.~903--916, 2018.
\bibitem{b6}
J.~Son et al., ``72-hour forecasting of hourly relativistic electron fluxes at
GEO by deep learning,'' \emph{Space Weather}, vol.~20, e2022SW003153, 2022.
\bibitem{b7}
H.~Zhang et al., ``A prediction model of relativistic electrons at
geostationary orbit using the EMD-LSTM network,''
\emph{Space Weather}, vol.~20, e2022SW003126, 2022.
\bibitem{b8}
X.~Sun et al., ``Enhancing \(\ge2\) MeV electron fluence predictions in GEO
through three deep learning strategies,'' \emph{Space Weather}, vol.~23,
e2024SW003975, 2025.
\bibitem{b9}
W.~Tan et al., ``Comparative analysis of TPA-LSTM and Transformer models
for forecasting GEO radiation belt electron fluxes,'' \emph{Space Weather}, vol.~22, e2023SW003714, 2024.
\bibitem{b10}
X.~Sun et al., ``A modeling study of \(\ge2\) MeV electron fluxes in GEO at
different prediction time scales based on LSTM and transformer networks,'' \emph{J. Space Weather Space Clim.},
vol.~14, p.~20, 2024.
\bibitem{b11}
E.~Camporeale, ``The challenge of machine learning in space weather,''
\emph{Space Weather}, vol.~17, pp.~1166--1207, 2019.
\bibitem{b12}
S.~K. Morley et al., ``Evaluating the performance of radiation belt models,''
\emph{Space Weather}, vol.~16, pp.~1521--1548, 2018.
\bibitem{b13}
A.~G. Ling et al., ``Identifying solar wind drivers of radiation belt
responses,'' \emph{Space Weather}, vol.~8, no.~10, 2010.
\bibitem{b14}
N.~Y.~Ganushkina et al., ``Forecasting GOES 15 \(>2\) MeV electron fluxes,''
\emph{Space Weather}, vol.~18, e2019SW002416, 2020.
\bibitem{b15}
S.~Sinha et al., ``PreMevE update,'' \emph{Space Weather}, vol.~19,
e2021SW002773, 2021.
\bibitem{b16}
A.~Vaswani et al., ``Attention is all you need,'' in \emph{Proc.\ Adv. Neural Inf. Process. Syst.\ (NeurIPS)}, 2017,
pp.~5998--6008.
\end{thebibliography}
\end{document}